% ステップ0 検証用: 色なし版（学習時の入力に相当）
\documentclass[uplatex,dvipdfmx]{jsarticle}

\begin{document}

\section{はじめに}

本研究では、ブロックエディタで編集した内容をコンパイル可能なTeXコードとして
出力するツールを開発する。本節ではその背景と目的について述べる。
既存のツールでは、TeXの文法を習得していない利用者が文書を作成することが難しい。

\subsection{目的}

PDFを読み込み、自作Transformerによるブロック分類を行い、
ブロックエディタ上で編集可能な形に変換することを目的とする。

\begin{figure}[h]
\centering
\fbox{\rule{0pt}{3cm}\rule{5cm}{0pt}}
\caption{システム全体の構成}
\end{figure}

\begin{table}[h]
\centering
\caption{対応ブロック種別}
\begin{tabular}{|l|l|}
\hline
種別 & 対応フェーズ \\
\hline
見出し & Phase 1 \\
段落 & Phase 1 \\
表 & Phase 2 \\
\hline
\end{tabular}
\end{table}

\end{document}